\BeleskeID{09_2-s003}
% element: 09_2-s003:e04
Visok napon na drejnu pogonskog tranzistora omogućava zasićenje čim se formira kanal. Za $T_D$ uslov je $V_{DS,D}\geq V_{GS,D}-V_{Tn,D}$, odnosno $V_O\geq V_I-V_{Tn,D}$. Opterećenje i dalje radi u omskoj oblasti. Polazni izraz za ravnotežu struja glasi
\[\frac{k_{n,L}}2\left(2V_{DS,L}(V_{GS,L}-V_{Tn,L})-V_{DS,L}^2\right)=\frac{k_{n,D}}2(V_{GS,D}-V_{Tn,D})^2\]
Zamenjuju se $V_{GS,L}=0$, $V_{DS,L}=V_{DD}-V_O$, $V_{GS,D}=V_I$ i $V_{DS,D}=V_O$. Množenjem sa dva i izvlačenjem faktora $V_{DD}-V_O$ dobija se zavisnost na slajdu.

Posle zamene napona ravnoteža glasi
\[\frac{k_{n,L}}2\left(2(V_{DD}-V_O)(-V_{Tn,L})-(V_{DD}-V_O)^2\right)=\frac{k_{n,D}}2(V_I-V_{Tn,D})^2\]
